% Research lemma; deliberately unregistered pending parent integration decision.
\paragraph{An intersection identity for the degree-unit SNC gluing term.}
Let $Z$ be regular of dimension $d+1$, with $d\ge2$, and let
$E=\bigcup_iE_i$ be a projective simple normal crossing divisor.
Let $D=\sum_i a_iE_i$, where all $a_i>0$ are rational.
Choose an ample line bundle $H$ on $E$. Suppose that for every
irreducible double stratum $C_{ij}\subset E_i\cap E_j$,
\[
                  \lambda_{ij}
             =\deg\bigl((-D)|_{C_{ij}}\,H^{d-2}\bigr)>0.
\]
This holds, in particular, if $\mathcal O_Z(-D)|_E$ is ample after
clearing the denominators of $D$.
Choose an orientation for each double stratum. Then the restriction-degree
map
\[
 \bigoplus_i N_1(E_i)_{\mathbb Q}
       \longrightarrow\bigoplus_{\text{double strata}}\mathbb Q,
 \qquad (z_i)_i\longmapsto
       \bigl(\deg(z_i\cdot C_{ij})-
                         \deg(z_j\cdot C_{ij})\bigr)_{i\to j}
\]
is surjective. Here each double stratum is a divisor on both adjacent
components; if an intersection is disconnected, we index its components
separately.

\begin{proof}
It suffices to prove that a real functional annihilating the image is
zero. Write its coefficients antisymmetrically as $t_{ij}=-t_{ji}$,
with the same convention on every connected double component.
Annihilation means that on each $E_i$ the numerical divisor class
\[
                         R_i=\sum_jt_{ij}[C_{ij}]
\]
is zero. Put $u_{ij}=t_{ij}/(a_i a_j)$. Thus
$\sum_ja_ju_{ij}[C_{ij}]=R_i/a_i$ is numerically zero.
Square these divisors, intersect with $H^{d-2}$, multiply the $i$th
identity by $a_i$, and sum. A double stratum contributes
\[
 a_i a_j u_{ij}^2\,
 \deg\bigl((a_jN_{C_{ij}/E_i}
                    +a_iN_{C_{ij}/E_j})H^{d-2}\bigr).
\]
By the transverse Cartier equations,
$N_{E_i/Z}|_{C_{ij}}=N_{C_{ij}/E_j}$ and conversely.
Consequently the displayed degree is the contribution of $D|_{C_{ij}}$,
minus the contributions from all third components meeting that stratum.

At a triple stratum $T\subset E_i\cap E_j\cap E_k$, put
$n_T=\deg(H^{d-2}|_T)>0$ and use cyclic orientations $ij,jk,ki$.
The three subtracted diagonal terms and the three cross terms give
exactly
\[
             -a_i a_j a_k n_T
                       (u_{ij}+u_{jk}+u_{ki})^2.
\]
There are no additional terms: all intersections of two boundary
divisors on a fixed $E_i$ are the indicated transverse triple strata.
We have therefore proved the exact identity
\begin{equation*}
 0=-\sum_{\text{double strata}}a_i a_j\lambda_{ij}u_{ij}^2
   -\sum_{\text{triple strata}}a_i a_j a_k n_T
                          (u_{ij}+u_{jk}+u_{ki})^2.
\end{equation*}
Every coefficient in the first sum is positive and every coefficient
in the second sum is nonnegative. Thus all $u_{ij}$ and all $t_{ij}$
vanish. This proves surjectivity.
\end{proof}

After tensoring the map with $k^\times\otimes\mathbb Q$, the same
surjectivity says that all degree-unit edge vectors can be obtained
by restricting curve classes multiplied by units from the components.
For a curve class on $E_i$, use its rational Chern-character class in
$K_0(E_i)^{(d-1)}$ and multiply by the chosen unit in $K_1(k)^{(1)}$;
restriction to a double stratum and degree give precisely the map
above. Thus a proposed top $K_1$ gluing obstruction based only on
these edge degree-unit vectors vanishes under the hypotheses.

This is not a theorem about the entire ordinary top Chow quotient.
Higher-stratum $K$-groups, the image of $K_0(Z)$ inside $K_0(E)$,
and realization of ordinary relations by actual families on $Z$
are separate questions. In particular, preserving a dual-complex
loop does not preserve this degree-unit gluing quotient.
